\section*{अध्याय \ref{ch:b1:optical-mineralogy} --- प्रकाशिकी खनिजविज्ञान}

\begin{solution}{exo:b1:optical-mineralogy:1}
स्फटिक: $\qty{30000}{nm}\times0.009 = \qty{270}{nm}$, प्रथम-कोटि का सफ़ेद।
कैल्साइट: $\qty{30000}{nm}\times0.172 = \qty{5160}{nm}$, चार्ट से बहुत परे: एक हल्का उच्च-कोटि का सफ़ेद, हल्के
गुलाबी या हरे आभास के साथ।
\end{solution}

\begin{solution}{exo:b1:optical-mineralogy:2}
$I/I_{\max} = \sin^2 2\theta$: \ang{0} पर $0$; \ang{22.5} पर $\sin^2\ang{45} = 0.5$; \ang{45} पर $1$; \ang{67.5} पर $0.5$;
\ang{90} पर $0$।
\end{solution}

\begin{solution}{exo:b1:optical-mineralogy:3}
\omterm{def:b1:rock-forming-minerals:groups}{गार्नेट} घनीय है, अतः समदैशिक: प्रकाश उसे उसी कंपन में छोड़ता है जिसमें उसने प्रवेश किया, ध्रुवक के अनुदिश, और विश्लेषक
उसे रोक देता है; हर स्थिति में $\Delta = 0$। अपने \omterm{def:b1:optical-mineralogy:optic-axis}{प्रकाशिक अक्ष} के लंबवत काटा गया स्फटिक भी उस अक्ष के अनुदिश चलते
प्रकाश के लिए $n_1 = n_2$ रखता है, अतः $\Delta = 0$: वह भी काला रहता है। \omterm{def:b1:field-lab-methods-1:microscope}{समतल-ध्रुवित प्रकाश} में \omterm{def:b1:rock-forming-minerals:groups}{गार्नेट} का \omterm{def:b1:optical-mineralogy:relief}{उभार} बहुत ऊँचा
है ($1.83$, $1.54$ के सामने), स्फटिक का लगभग नहीं।
\end{solution}

\begin{solution}{exo:b1:optical-mineralogy:4}
$|n - 1.54|$: अल्माँडीन $0.29$, फ़ॉर्स्टेराइट $0.11$, ऑर्थोक्लेज़ $0.02$, स्फटिक $0.01$; \omterm{def:b1:optical-mineralogy:relief}{उभार} उसी क्रम में, ऊँचे से निचले तक।
ऑर्थोक्लेज़ का अपवर्तनांक नीचा है: ऋणात्मक \omterm{def:b1:optical-mineralogy:relief}{उभार}, और फ़ोकस ऊपर उठाने पर उसकी \omterm{def:b1:optical-mineralogy:relief}{बेक रेखा} कण से बाहर, आरोपण में खिसकती है।
\end{solution}

\begin{solution}{exo:b1:optical-mineralogy:5}
$\delta = 650/30000 = 0.022$। \ang{56} और \ang{124} पर विदलनों वाला और लगभग $0.020$ के \omterm{def:b1:optical-mineralogy:birefringence}{द्विअपवर्तन} वाला एक हरा \omterm{def:b1:optical-mineralogy:pleochroism}{बहुवर्णी}
\omterm{def:b1:rock-forming-minerals:mineral}{खनिज}: हॉर्नब्लेंड।
\end{solution}

\begin{solution}{exo:b1:optical-mineralogy:6}
$\tan\ang{30} = 0.5774$। स्फटिक: $(0.013/1.54)\times0.5774 = \qty{0.00487}{rad}
= \ang{0.28}$। फ़ॉर्स्टेराइट: $(0.111/1.54)\times0.5774 = \qty{0.0416}{rad} =
\ang{2.38}$। अनुपात $0.111/0.013 = 8.5$। ठीक-ठीक: $\sin r = 1.54\sin\ang{30}/1.651
= 0.4664$, $r = \ang{27.80}$, विचलन \ang{2.20}: इतने बड़े अपवर्तनांक अंतर के लिए प्रथम-कोटि सूत्र \qty{8}{\percent} अधिक
आँकता है।
\end{solution}

\begin{solution}{exo:b1:optical-mineralogy:7}
$270 + 540 = \qty{810}{nm}$, द्वितीय-कोटि का हरा, जब कण और पट्टी की धीमी दिशाएँ समांतर हों (विलंब जुड़ते हैं); $|270 - 540| =
\qty{270}{nm}$, फिर प्रथम-कोटि का सफ़ेद, जब कण की धीमी दिशा पट्टी की तेज़ दिशा के समांतर हो।
\end{solution}

\begin{solution}{exo:b1:optical-mineralogy:8}
पूर्ण मोटाई पर $\Delta = \qty{30000}{nm}\times0.052 = \qty{1560}{nm}$। किनारे ($\Delta = 0$) से अंदर की ओर \omterm{def:b1:optical-mineralogy:retardation}{मंदन} 542 और
\qty{1074}{nm} पार करता है पर \qty{1605}{nm} नहीं: दो लाल-बैंगनी झालरें।
\end{solution}

\begin{solution}{exo:b1:optical-mineralogy:9}
किसी कण पर मापा गया \omterm{def:b1:optical-mineralogy:extinction}{विलोपन कोण} इस पर निर्भर करता है कि कण कैसे काटा गया है: केवल एक विशेष अभिविन्यास की काटें सबसे बड़ा
कोण दिखाती हैं, और अन्य शून्य तक कम। सबसे बड़ा मान, लगभग \ang{37}, वह है जिससे तुलना करनी है: वह डायोप्साइड
(लगभग \ang{38}) से मेल खाता है, हॉर्नब्लेंड (अधिकतम \ang{19}) से नहीं।
\end{solution}

\begin{solution}{exo:b1:optical-mineralogy:10}
विश्लेषक को $x$ के अनुदिश रखने पर, दो तरंगों के $x$ घटक $E_0\cos\theta\cdot\cos\theta$ और
$-E_0\sin\theta\,e^{-i\varphi}\cdot(-\sin\theta)$ हैं, अतः $A = E_0(\cos^2\theta + \sin^2\theta\,e^{-i\varphi})$ और
\[
  |A|^2 = E_0^2\bigl(\cos^4\theta + \sin^4\theta + 2\sin^2\theta\cos^2\theta\cos\varphi\bigr)
  = E_0^2\bigl(1 - 2\sin^2\theta\cos^2\theta\,(1 - \cos\varphi)\bigr),
\]
$\cos^4\theta + \sin^4\theta = 1 - 2\sin^2\theta\cos^2\theta$ का उपयोग करके। चूँकि
$2\sin^2\theta\cos^2\theta = \frac12\sin^2 2\theta$ और $1 - \cos\varphi =
2\sin^2(\varphi/2)$, $I_\parallel = I_0(1 - \sin^2 2\theta\sin^2(\pi\Delta/\lambda))$।
हर तरंगदैर्घ्य पर $I_\parallel + I_\perp = I_0$: समांतर ध्रुवकों के बीच के रंग \omterm{def:b1:field-lab-methods-1:microscope}{आड़े ध्रुवकों} के बीच के रंगों के पूरक हैं। \qty{550}{nm}
का एक कण, जो \omterm{def:b1:field-lab-methods-1:microscope}{आड़े ध्रुवकों} के बीच मैजेंटा है, समांतर ध्रुवकों के बीच हल्का पीला-हरा है।
\end{solution}

\begin{solution}{exo:b1:optical-mineralogy:11}
(क) $\Delta/\lambda$ $5000/700 = 7.14$ से $5000/400 = 12.5$ तक चलता है; वह पूर्ण संख्याएँ 8 से 12 पार करता है: पाँच शून्य।
(ख) शून्य $\lambda_k = \Delta/k$ पर हैं; दो क्रमागत $\Delta/k - \Delta/(k + 1) =
\Delta/(k(k + 1)) \approx \lambda^2/\Delta$ दूर हैं, यानी \qty{550}{nm} के निकट $550^2/5000 =
\qty{60.5}{nm}$। (ग) वर्णक्रम में फैली पाँच संकरी काली पट्टियाँ हर रंग का थोड़ा हटाती हैं, और शेष औसतन आधे भार के साथ
गुज़रता है: आँख एक हल्का, लगभग तटस्थ आभास देखती है।
\end{solution}

\begin{solution}{exo:b1:optical-mineralogy:12}
$\sin^2 2\theta$ का एक चक्कर पर माध्य $\frac1{2\pi}\int_0^{2\pi}\sin^2 2\theta\,
\mathrm d\theta = \frac12$ है, अतः माध्य तीव्रता
$\frac12 I_0\sin^2(\pi\Delta/\lambda)$ है। वह $\delta = \Delta/30000$ परिकलित करती है, बजाय $\Delta/25000$ के (nm में): उसके सभी
मान सही मानों के $25/30 = 0.83$ गुना हैं, \qty{17}{\percent} कम।
\end{solution}

\begin{solution}{pb:b1:optical-mineralogy:1}
\textbf{1.} B: गहरी किनारियों का अर्थ उच्च \omterm{def:b1:optical-mineralogy:relief}{उभार} है; A की मुश्किल से दिखती सीमाओं का अर्थ निम्न \omterm{def:b1:optical-mineralogy:relief}{उभार}।

\textbf{2.} A: $(1.544 + 1.553)/2 = 1.5485$, $|n - n_m| = 0.0085$। B:
$(1.664 + 1.672 + 1.694)/3 = 1.6767$, $|n - n_m| = 0.1367$।

\textbf{3.} विचलन $|n - n_m|$ के अनुपात में हैं: $0.1367/0.0085 =
16.1$।

\textbf{4.} B का अपवर्तनांक आरोपण से ऊँचा है: धनात्मक \omterm{def:b1:optical-mineralogy:relief}{उभार}।

\textbf{5.} \omterm{def:b1:optical-mineralogy:pleochroism}{बहुवर्णता}; बायोटाइट।

\textbf{6.} \ang{56} और \ang{124} पर दो विदलन: एक \omterm{def:b1:rock-forming-minerals:groups}{एम्फ़ीबोल}।

\textbf{7.} $\qty{30000}{nm}\times0.009 = \qty{270}{nm}$: प्रथम-कोटि का सफ़ेद।

\textbf{8.} हर कण अपने अभिविन्यास में काटा गया है और $|n_1 - n_2|$ का, 0 और $\delta$ के बीच, एक मान दिखाता है; केवल \omterm{def:b1:optical-mineralogy:optic-axis}{प्रकाशिक
अक्ष} के समांतर काटे गए कण स्वयं $\delta$ दिखाते हैं। कई कणों में सबसे ऊँचा रंग वह है जिसका उपयोग करना है।

\textbf{9.} $d = 342/0.009 = \qty{38000}{nm} = \qty{38}{\micro m}$।

\textbf{10.} $(38 - 30)/30 = \qty{26.7}{\percent}$।

\textbf{11.} $\sin^2(\pi\times342/550) = 0.861$; \qty{30}{\micro m} पर,
$\sin^2(\pi\times270/550) = 0.999$।

\textbf{12.} $\qty{38000}{nm}\times0.008 = \qty{304}{nm}$ से
$\qty{38000}{nm}\times0.013 = \qty{494}{nm}$ तक: सफ़ेद से नारंगी-लाल तक, सभी \omterm{def:b1:optical-mineralogy:interference-colour}{प्रथम कोटि} के अंत \qty{542}{nm} से नीचे।

\textbf{13.} 1074 और \qty{1605}{nm} के बीच: \omterm{def:b1:optical-mineralogy:interference-colour}{तृतीय कोटि}।

\textbf{14.} $\delta = 1140/38000 = 0.030$।

\textbf{15.} $1140/30000 = 0.038$, जो एक \omterm{def:b1:rock-forming-minerals:groups}{ओलिवीन} सुझाएगा। एक मैग्नीशियम-ओलिवीन और स्फटिक मिलकर \omterm{def:b1:rock-forming-minerals:groups}{पाइरॉक्सीन} देते
हैं, \ce{Mg2SiO4 + SiO2 -> 2 MgSiO3}; स्फटिक के साथ क्रिस्टलित चट्टान दोनों नहीं रख सकती।

\textbf{16.} $0.030$ डायोप्साइड का \omterm{def:b1:optical-mineralogy:birefringence}{द्विअपवर्तन} है; उसका ऊँचा \omterm{def:b1:optical-mineralogy:relief}{उभार}, उसका रंगहीन होना और लगभग \ang{87} पर उसके दो विदलन मेल
खाते हैं: B एक क्लाइनोपाइरॉक्सीन, डायोप्साइड है।

\textbf{17.} $\qty{30000}{nm}\times0.030 = \qty{900}{nm}$: द्वितीय-कोटि का पीला।

\textbf{18.} \qty{8}{\micro m} तक; तब स्फटिक अधिक-से-अधिक \qty{270}{nm} दिखाता है, प्रथम-कोटि का सफ़ेद।

\textbf{19.} \omterm{def:b1:optical-mineralogy:twin}{बहुसंश्लिष्ट यमलन}: D \omterm{def:b1:rock-forming-minerals:groups}{प्लैजियोक्लेज़} है।

\textbf{20.} \omterm{def:b1:rock-forming-minerals:groups}{प्लैजियोक्लेज़} का संघटन (उसका एनॉर्थाइट अंश)।

\textbf{21.} \ang{17} के \omterm{def:b1:optical-mineralogy:extinction}{विलोपन कोण} वाला एक \omterm{def:b1:rock-forming-minerals:groups}{एम्फ़ीबोल}: हॉर्नब्लेंड (\ang{15}--\ang{19}), \omterm{def:b1:rock-forming-minerals:groups}{पाइरॉक्सीन} नहीं (डायोप्साइड के लिए \ang{38})।

\textbf{22.} $\qty{38000}{nm}\times0.020 = \qty{760}{nm}$: द्वितीय-कोटि का हरा।

\textbf{23.} यह \omterm{def:b1:optical-mineralogy:optic-axis}{प्रकाशिक अक्ष} के लंबवत काटा गया है: उसके दो अपवर्तनांक उस अक्ष के अनुदिश चलते प्रकाश के लिए बराबर हैं, $\Delta = 0$,
और वह एक समदैशिक कण की तरह काला रहता है।

\textbf{24.} A स्फटिक, B डायोप्साइड, C बायोटाइट, D \omterm{def:b1:rock-forming-minerals:groups}{प्लैजियोक्लेज़}, E हॉर्नब्लेंड: कुछ स्फटिक वाली एक मध्यवर्ती आग्नेय चट्टान,
जैसे एक स्फटिक डायोराइट।

\textbf{25.} काट \textbf{\qty[text-series-to-math=true]{38}{\micro m}} मोटी है।
\end{solution}
